\section{有限型の表示族と三捩れ因子を越える楕円境界}
\label{sec:algebraic-boundary}
形式的な圏の認識と、代数的な表示族の存在は区別する。
明示中心式のおかげで、後者については有限個の開条件へ進める。

\begin{proposition}[有限型底上の全次数平坦な表示]\label{prop:algebraic-presentation-flat}
定理\ref{thm:eight-formal-central-lift}の八変数有理係数表示は、
選んだ点を含む有限型smoothな八次元底のZariski開部分 $U$ 上に定義できる。
底はまず
\[
 \eta\zeta p(p^3+1)\ne0,\qquad p=1+\kappa\lambda
\]
に制限し、中心式の有限個の分母と、$N_{x^3}/\lambda^2$ をさらに可逆化する。
この $U$ 上で包絡と中央三次商は次数別有限局所自由で、任意の底変更に適合する。
各幾何学的fiberの包絡はnoetherian Koszul CY AS4の整域で、三次元は正則中央である。
この命題は全fiberのH3 admissibilityを結論しない。
\end{proposition}
\begin{proof}
中心式は有限個の有理係数を持ち、共通 $\lambda^2$ の消去は局所環の厳密な所属から得た。
従って形式冪級数の無限個の係数を代数化する必要はなく、有限個の分母を除く開部分で定義される。
残る $N_{x^3}/\lambda^2$ の $\lambda=0$ 値は
$9\eta^2\zeta^2(1+q_0)$ なので、必要な開部分は選んだ点を含む。

斉次先頭potentialを $X=x,Y=\eta y,Z=\zeta z$ で書くと、
三次元Sklyaninの三係数は
\[
 (p,-q_0,3\lambda\eta\zeta)
\]
となる。最初の二係数は非零で、その立方は $p^3$ と $-1$ で異なる。
\cite[Theorem1.1, Corollary6.1]{IS}により、全fiberでKoszul・polynomial Hilbert・CY3である。
同一次本文の仮定を今回確認した。
標準の三次元AS正則代数のnoetherian・domain定理も、従来と同じ既知入力として使う。
これらの性質をISのHilbert計算だけから推論してはいない。
\cite[Theorem3.3.2(i)]{EG}で任意の二次・一次potential項を加えたfiltered PBW変形を取り、
そのRees環と正則中央一次拡大の認識 \cite{LBS} により、fiberの包絡は主張したAS4となる。

全次数での\emph{底上の}flat性も説明する。
次数 $n$ の表示行列は有限行列で、全fiberの商次元は $\binom{n+3}{3}$ で一定である。
底 $U$ はreduced Jacobson環上の開部分なので、rankより一つ大きい小行列式は恒等的に零であり、
rankを与える小行列式は $U$ を覆う。
従って各次数の商は有限局所自由である。この議論に次数ごと異なる開部分を追加していない。

中心恒等式は稠密開上で成立し、包絡の各次数は $U$ 上torsion-freeなので $U$ 全体へ延長する。
既存の正規形の分母を除いたため、$x^3$ 係数1の三次元は各fiberで非零である。
fiberの包絡は整域なので正則である。
各次数の中央三次乗法は有限局所自由加群間でfiberごとに単射となり、
その余核も有限局所自由である。これが商のflat性と任意底変更を与える。
\end{proof}

\begin{proposition}[正規化した楕円曲線の明示式]\label{prop:algebraic-boundary-cubic}
\[
 C=\kappa+\kappa^2\lambda+
       \left(\frac{\kappa^3}{3}+9\eta^3\zeta^3\right)\lambda^2,
 \qquad
 F=q_0p(x^3+\eta^3y^3+\zeta^3z^3)+Cxyz
\]
と置く。
斉次先頭環のmultilinear関係行列
\[
 M=\begin{pmatrix}
 3\lambda x&-q_0z&py\\
 pz&3\lambda\eta^3y&-q_0x\\
 -q_0y&px&3\lambda\zeta^3z
 \end{pmatrix}
\]
は
\[
 \det M=3\lambda F
\]
を満たす。
従って $\lambda\ne0$ の点曲線の、$\lambda=0$ を越える正規化した三次族は $F=0$ である。
\[
 \Delta_E=C^3+27p^3\eta^3\zeta^3\ne0
\]
をさらに課せばこれはsmoothな相対楕円曲線となる。
選んだ点で $\Delta_E=170$ であり、この開条件は形式近傍を失わない。
$\lambda=0$ では関係行列は射影平面全体でrank2で、対応する写像は
\[
 (x:y:z)\longmapsto(x:q_0y:q_0^2z)
\]
である。選択した三次元の多重線型化は、$F|_{\lambda=0}$ の非零定数倍となる。
\end{proposition}
\begin{proof}
行列式と $3\lambda C=p^3-1+27\lambda^3\eta^3\zeta^3$ は直接展開である。
三つの立方項の係数が非零のHesse三次式について、三偏微分の積を使えば
特異点が存在する条件は $C^3+27p^3\eta^3\zeta^3=0$ と分かる。
$\lambda=0$ でこの式は $\kappa^3+27\eta^3\zeta^3$、指定点では $8+27\cdot2\cdot3=170$ となる。
\texttt{verify\_algebraic\_boundary.py}はこれらを全係数で検証する。
同検証器は $\lambda=0$ の各射影patchで全二次小行列式の零点がないこともGroebner基底 $(1)$ により確認する。
表示した対角写像を $M$ に代入すれば零となる。
三次元の多重線型評価では $xyz$ は $q_0^2xyz$、各三乗はそのままなので、
$\kappa q_0^2xyz+x^3+\eta^3y^3+\zeta^3z^3=q_0^{-1}F|_{\lambda=0}$ である。
この点曲線の明示化だけを、全相対THCR認識や全H3開性と同一視しない。
\end{proof}

\begin{theorem}[相対THCR境界と全fiberの整域性]\label{thm:relative-THCR-boundary}
前命題の $U$ を、$\Delta_E\ne0$ と、必要ならさらに選んだ点を含む有限の開部分へ制限する。
このとき $E=(F=0)\to U$ には関係行列から定まる相対自己同型 $\sigma$ があり、
$L=\OO_E(1)$ とすると
\[
 R_U/(t)\simeq B(E,L,\sigma)
\]
が任意の底変更に適合して成立する。
全幾何学的fiberで $t$ は $R$ 上正則中央であり、$R$ はnoetherian整域である。
この結論は全fiberのH3 smoothnessを含まない。
\end{theorem}
\begin{proof}
まず関係式を右側の点の座標に関してまとめると、もう一つの行列
\[
 M^{\rm r}=\begin{pmatrix}
 3\lambda x&pz&-q_0y\\
 -q_0z&3\lambda\eta^3y&px\\
 py&-q_0x&3\lambda\zeta^3z
 \end{pmatrix}
\]
を得て、その行列式も $3\lambda F$ である。
$\lambda=0$ では左右とも射影平面全体でrank2であり、左右の核は互いに逆の対角写像を与える。
$E\to U$ はprojectiveなので、左右のrank1以下の閉部分の像は閉である。
指定fiberにその像はないため、それらを除く有限の開部分は指定点を含む。

この開部分で $M$ の核は直線束となり、$\sigma:E\to\PP^2_U$ を与える。
$\lambda\ne0$ では $M^{\rm r}$ の行列式から像は $E$ に入り、
$E$ のflat性と $\lambda\ne0$ の稠密性からこの性質は全体へ延長する。
$M^{\rm r}$ から得る逆写像も同様で、核の一意性により互いに逆である。
$\lambda=0$ での式は既に検証した対角写像である。

この相対データからTHCRを作る。
各正次数のline bundle $L_n=L\otimes\sigma^*L\otimes\cdots\otimes(\sigma^{n-1})^*L$ はfiber次数 $3n$ であり、
$H^1$ が消えるので、その切断はrank $3n$ の有限局所自由加群で底変更に適合する。
THCRは次数1で生成される。fiber上での説明を与えると、$L$ のbase-point-freeな二切断のpencilを取り、
その完全列を他の正次数line bundleでtensorすれば、次数3より大きい場合の乗法全射を得る。
二つの次数3line bundleが非同型なら同じpencilの議論でよく、同型なら
平面三次曲線が二次関係を持たないことから乗法全射を得る。
従って生成性は全fiberで成立し、Nakayamaにより相対的にも成立する。
関係行列が $E$ と $\sigma$ のgraph上で消えるため、$S_U=B_U/(t)$ からTHCRへの自然全射がある。

中央三次元 $\overline G$ がこの全射で零になることを確認する。
まずgenericな $\sigma^3$ は恒等でない。
これは有理section $(0:\zeta:-\eta)$ の三反復を用いて厳密に検証できる。
\texttt{verify\_boundary\_graph.py}は
$\lambda=1/2,\kappa=2,\eta=\zeta=1$ でその軌道を計算し、
三反復後の第一座標が非零であることを確認する。
これはgenericな有理写像の恒等式を一つの正確な反例で否定しているのであり、全fiberを標本から外挿していない。

THCRは、generic曲線の関数体 $K(E)$ のskew Laurent環へ埋め込める。
次数1の切断の比は $K(E)$ を生成するため、非零な中央次数3元があれば $\sigma^3$ は関数体上恒等となる。
genericにはそうでないので、中央元 $\overline G$ のTHCRへの像は零である。
この像は $U$ 上の有限局所自由な三次切断加群のsectionなので、稠密generic点で零なら全体で零である。
従って全射は $S_U/(\overline G)$ を通る。
全fiberで $S$ はAS3整域、$\overline G\ne0$ なので、商のHilbert値は $1,3,6,9,\ldots$。
THCRのHilbert値と等しく、両辺は次数別有限局所自由なので、全射は各次数で同型である。

$t,G$ は各fiberの包絡上で正則中央列である。
実際 $t$ は包絡上正則であり、$S$ 上の $\overline G$ も正則である。
このため順序を入れ替えて $R$ 上の $t$ も正則となる。
商 $R/(t)$ はsmooth楕円曲線のTHCRなので整域であり、
connected gradingによる有限の $t$-進次数を取れば $R$ 自体も整域となる。
noetherian性は既に得たnoetherian包絡の商であることから従う。
\end{proof}

\begin{corollary}[幾何学的一般点のH3性]\label{cor:eight-generic-H3}
この明示八変数族の幾何学的一般点はH3である。
同時に、選んだ点での全Artin形式分類は定理\ref{thm:intrinsic-formal-eight}で与えられる。
これら二つの事実だけから、族の全H3 locusがZariski開であるとは結論しない。
\end{corollary}
\begin{proof}
$\lambda,\kappa,\eta,\zeta,a,b,c$ の関数体を $K$ とし、$D$ を独立な超越元とする。
非斉次包絡 $V_K$ はCY3 PBW環で有限大域次元を持ち、中心三次元 $H$ は $D$ に依存しない。
その最高記号は正次数の非零元なので、$H$ は $K$ 上超越的である。
\[
 V_K\otimes_KK(D)/(H+D)
 \simeq V_K\otimes_{K[H]}K(H)
\]
は中心局所化であり、有限大域次元を持つ。
前定理でsmoothな楕円THCR境界を認識したので、既存の補題\ref{lem:generic-D}と
入力\ref{ass:recognition}(3)を適用できる。
必要な分離的体拡大と幾何学的fiberの認識は同じ入力に従う。
これは独立な中心値 $D$ のgeneric点についての証明である。
任意の数値 $D$ や有限位数を横断する一様なH3開性を主張する議論ではない。
\end{proof}
